\begin{singlespace}

\indent Penyakit jantung koroner kronis (CCD) berkontribusi sebagai beban terbesar pada penyakit
kardiovaskular secara global, namun pelaksanaan sesuai panduan klinis pada praktik rutin masih rendah:
hanya 1,0\% dari 4.548 pasien pada studi INTERASPIRE yang mencapai pelaksanaan yang optimal.
Sistem rekomendasi klinis dapat mengatasinya, tetapi harus dilakukan penyatuan rekam
medis, panduan klinis, dan literatur biomedis secara sekaligus ke dalam satu struktur yang dapat dikueri,
menjembatani terminologi klinis yang heterogen, serta menghasilkan keluaran dalam bahasa alami
yang \emph{faithfulness}-nya dapat diverifikasi. Penelitian ini mengembangkan sistem
rekomendasi klinis CCD yang bersifat FHIR-native dan berbasis \emph{knowledge graph}
untuk menjawab ketiganya.

Rekam medis sintetis sejumlah 5.490 pasien dibangkitkan menggunakan Synthea, diekspor sebagai
\emph{bundle} HL7 FHIR~R4, dimuat ke basis data graf Neo4j, lalu disatukan dengan panduan
AHA/ACC CCD 2023 dan 4.801 abstrak PubMed menggunakan lapisan konsep bersama berbasis UMLS. Karena
observasi numerik ikut dimuat sebagai data graf, rekomendasi yang disyaratkan pada nilai
terukur---misalnya fraksi ejeksi ventrikel kiri atau indeks massa tubuh---dapat digerbangi
menurut nilai pasien yang sebenarnya, yang tidak dapat dilakukan jika berbasis kondisi penyakitnya
saja. Sistem juga menghasilkan \emph{grounded briefing} yang faktanya ditelusuri secara
deterministik dari graf dan model bahasa besar hanya merumuskan kalimatnya saja, lalu
diperiksa secara otomatis. Evaluasi yang mencakup 876 pasien CCD dilakukan terhadap standar rujukan
otomatis yang independen dari sistem.

Graf yang terbentuk memuat 1.804.379 simpul dan 8.796.425 relasi. Ablasi menunjukkan
\emph{observation gate} dapat menaikkan $F_1$ pada keselarasan terhadap panduan klinis dari 0,743 menjadi 0,810 melalui
kenaikan presisi dari 0,612 menjadi 0,684 dengan \emph{recall} yang praktis tidak berubah dan
seluruh perolehan itu terpusat pada 15 dari 68 rekomendasi yang memang bergerbang. Standar
rujukan otomatisnya sepakat dengan pelabelan manual terhadap 30 pasien pada 99,5\% sel dengan
Cohen's $\kappa = 0{,}991$ (interval kepercayaan 95\% 0,984---0,996). Seluruh 30 \emph{briefing} yang dibangkitkan lolos pemeriksaan
\emph{faithfulness} otomatis, tanpa satu pun rekomendasi maupun sitasi yang dikarang, poin
wajib yang terlewat, ataupun klaim antikoagulan yang berlebih. Telaah oleh dokter spesialis
jantung menilai 77,7\% dari 247 rekomendasi bagi 15 pasien tepat, dengan penilaian belum
tepat terpusat pada terapi antitrombotik.

Pemuatan observasi numerik sebagai data kelas satu karenanya menaikkan presisi keselarasan
terhadap panduan klinis secara terukur. FHIR menyediakan observasi itu secara paling
langsung, tetapi ekspor tabular pun dapat mereproduksi gerbang yang sama. Penelusuran yang
deterministik menjadikan \emph{faithfulness} pada tingkat sitasi dapat diputuskan secara
mekanis, bukan aproksimatif seperti pada \emph{retrieval augmented generation} (RAG), tetapi
tidak menjamin bahwa setiap kalimat setia secara makna. Keterbatasannya pun
dapat diukur: ekspor FHIR tidak memuat tanggal berhenti pengobatan sehingga logika panduan yang
bergantung pada durasi hanya sebagian yang dapat diterapkan pada saat penelusuran dan seluruh angkanya berasal dari
rekam medis sintetis.

\vspace{6pt}
\noindent{Kata kunci} : sistem rekomendasi klinis, \emph{knowledge graph}, HL7 FHIR,
penyakit jantung koroner kronis, keselarasan panduan

\end{singlespace}
